\documentclass[11pt,a4paper]{article}

% ----------------------------------------------------------------------
%  PACKAGES
% ----------------------------------------------------------------------
\usepackage[utf8]{inputenc}
\usepackage[T1]{fontenc}
\usepackage[french]{babel}
\usepackage{lmodern}
\usepackage{microtype}
\usepackage{geometry}
\geometry{margin=2.2cm, top=2.4cm, bottom=2.4cm}

\usepackage{amsmath,amssymb}
\usepackage{xcolor}
\usepackage{tcolorbox}
\tcbuselibrary{skins,breakable}
\usepackage{enumitem}
\usepackage{booktabs}
\usepackage{array}
\usepackage{longtable}
\usepackage{fancyhdr}
\usepackage{titlesec}
\usepackage[hidelinks]{hyperref}

% ----------------------------------------------------------------------
%  CORRECTIONS D'ENCODAGE
% ----------------------------------------------------------------------
% 1) Hauteur d'en-tête (évite l'avertissement fancyhdr)
\setlength{\headheight}{15pt}

% 2) Prise en charge des caractères Unicode non couverts par T1
\usepackage{newunicodechar}

% Symbole « ∞ » du HTML
\newunicodechar{∞}{\ensuremath{\infty}}

% Symboles de la barre d'outils des lignes
\newunicodechar{✓}{\ensuremath{\checkmark}}          % valider
\newunicodechar{✎}{\ensuremath{\star}}               % modifier
\newunicodechar{▲}{\ensuremath{\blacktriangle}}      % monter
\newunicodechar{▼}{\ensuremath{\blacktriangledown}}  % descendre
\newunicodechar{✕}{\ensuremath{\times}}              % supprimer

% Symbole « ⌘ » (touche Cmd sur Mac)
\newunicodechar{⌘}{\textsf{Cmd}}

% Symboles mathématiques Unicode divers
\newunicodechar{≤}{\ensuremath{\leq}}
\newunicodechar{≥}{\ensuremath{\geq}}
\newunicodechar{≠}{\ensuremath{\neq}}

% ----------------------------------------------------------------------
%  COULEURS
% ----------------------------------------------------------------------
\definecolor{bleu}{HTML}{2563EB}
\definecolor{bleufonce}{HTML}{1E3A8A}
\definecolor{or}{HTML}{B45309}
\definecolor{orclair}{HTML}{FFFBEB}
\definecolor{vert}{HTML}{15803D}
\definecolor{rouge}{HTML}{B91C1C}
\definecolor{grisdoux}{HTML}{F1F5F9}

% ----------------------------------------------------------------------
%  STYLE
% ----------------------------------------------------------------------
\titleformat{\section}{\Large\bfseries\color{bleufonce}}{\thesection.}{0.6em}{}
\titleformat{\subsection}{\large\bfseries\color{bleu}}{\thesubsection}{0.5em}{}
\titleformat{\subsubsection}{\normalsize\bfseries\color{or}}{\thesubsubsection}{0.5em}{}

\setlist[itemize]{topsep=2pt,itemsep=2pt,parsep=0pt,leftmargin=1.4em}
\setlist[enumerate]{topsep=2pt,itemsep=2pt,parsep=0pt,leftmargin=1.6em}

\pagestyle{fancy}
\fancyhf{}
\fancyhead[L]{\small\color{bleufonce}\textbf{Mode d'emploi — Évaluation « Suites et limites »}}
\fancyhead[R]{\small\color{bleufonce}Term. gén. — Maths compl.}
\fancyfoot[C]{\small\thepage}
\renewcommand{\headrulewidth}{0.4pt}

% Boîte « À retenir »
\newtcolorbox{retenir}[1][À retenir]{
  colback=blue!5, colframe=blue!60!black, boxrule=0.8pt,
  arc=3pt, left=8pt, right=8pt, top=6pt, bottom=6pt,
  fonttitle=\bfseries\color{white}, coltitle=white,
  colbacktitle=bleufonce, title=#1, breakable
}

% Boîte « Attention »
\newtcolorbox{attention}[1][Attention]{%
  colback=red!4, colframe=red!60!black, boxrule=0.8pt,
  arc=3pt, left=8pt, right=8pt, top=6pt, bottom=6pt,
  fonttitle=\bfseries\color{white}, coltitle=white,
  colbacktitle=rouge, title=#1, breakable
}

% Boîte « Astuce »
\newtcolorbox{astuce}[1][Astuce]{%
  colback=green!4, colframe=green!45!black, boxrule=0.8pt,
  arc=3pt, left=8pt, right=8pt, top=6pt, bottom=6pt,
  fonttitle=\bfseries\color{white}, coltitle=white,
  colbacktitle=vert, title=#1, breakable
}

% ----------------------------------------------------------------------
%  TITRE
% ----------------------------------------------------------------------
\title{\vspace{-1.2cm}\textbf{\color{bleufonce}Mode d'emploi}\\[2pt]
       \large Application d'évaluation — \emph{Suites et limites}\\
       \normalsize Niveau : Terminale générale — option Mathématiques complémentaires}
\author{\textbf{Yann Merdy} \quad\texttt{\small yann.merdy@gmail.com}}
\date{\today}

% ======================================================================
\begin{document}
\maketitle
\thispagestyle{fancy}

\begin{retenir}[Résumé en cinq lignes]
\begin{enumerate}
  \item Je saisis mon \textbf{NOM} (en majuscules), mon \textbf{prénom}
        (initiale en majuscule), puis je coche éventuellement \emph{PAP}
        (tiers-temps) et je clique sur \textbf{Commencer le devoir}.
  \item Je rédige chaque question en \textbf{lignes de raisonnement} :
        \emph{une idée = une ligne}.
  \item La \textbf{première ligne} ne commence par aucun mot particulier ;
        chaque ligne \textbf{suivante} commence par \texttt{mais, ou, et, donc, or, ni, car}
        ou par \texttt{Finalement}.
  \item La \textbf{réponse} est la ligne qui commence par \texttt{Finalement}
        (elle est automatiquement encadrée en orange).
  \item Les mathématiques s'écrivent en \textbf{LaTeX direct, \emph{sans}
        taper les \texttt{\$}}. L'application les ajoute automatiquement :
        \verb|u_n = \dfrac{1}{n}| devient $u_n = \dfrac{1}{n}$.
        La formule n'est rendue qu'\textbf{après validation} de la ligne
        (\texttt{Ctrl+Entrée} ou bouton \textbf{✓}).
\end{enumerate}
\end{retenir}

\tableofcontents
\newpage

% ======================================================================
\section{Présentation générale}

Le document \texttt{TG-MATH-COMPLEMENTAIRES\_SUITES\_ET\_LIMITES.html} est une
\emph{application web autonome} (un seul fichier \texttt{.html}) qui permet de
composer et de rendre une évaluation de mathématiques sur les \textbf{suites et
les limites} en classe de Terminale générale, option Mathématiques complémentaires.

Le sujet est \textbf{généré aléatoirement} à chaque démarrage : les valeurs
numériques, les coefficients, les types de suites et les limites attendues
peuvent donc varier d'un élève à l'autre. La structure est toujours la même :

\begin{enumerate}
  \item \textbf{Exercice 1} — Formes indéterminées « classiques »
        (polynômes, quotients de polynômes).
  \item \textbf{Exercice 2} — Suites et fonctions trigonométriques
        (\(+\) encadrement, \(a\,n + b\sin n\), \dots).
  \item \textbf{Exercice 3} — Suites définies par des racines carrées
        (expression conjuguée).
  \item \textbf{Gros problème} — Étude complète d'une suite rationnelle
        (limite, écart \(v_n = u_n - \ell\), perturbation par \(\tfrac{\sin n}{n}\),
        produit \(n\,v_n\)).
\end{enumerate}

\subsection*{Ce que l'application fait pour vous}
\begin{itemize}
  \item Elle affiche l'énoncé et les questions, avec un rendu mathématique soigné
        (grâce à MathJax) ;
  \item elle \textbf{ajoute automatiquement les \texttt{\$}} autour des formules
        \LaTeX{} que vous tapez, puis les rend dans un aperçu dédié après
        validation de chaque ligne ;
  \item elle permet de \textbf{modifier, déplacer ou supprimer} n'importe quelle
        ligne de raisonnement (voir section~\ref{sec:edition}) ;
  \item elle chronomètre l'épreuve et affiche le temps restant ;
  \item elle \textbf{enregistre automatiquement} le travail toutes les 15 secondes
        dans le navigateur ;
  \item à la remise de la copie, elle produit un \textbf{fichier \texttt{.json}
        chiffré} à transmettre.
\end{itemize}

\subsection*{Ce que l'application \emph{ne} fait pas}
\begin{itemize}
  \item Elle ne corrige pas votre copie : elle compare seulement la
        \emph{dernière valeur numérique détectée} dans votre ligne
        \texttt{Finalement} à la réponse attendue (voir la
        section~\ref{sec:detection}) ;
  \item elle ne vous \emph{empêche pas} d'écrire une erreur : c'est votre
        raisonnement qui est évalué par votre professeur, pas seulement la
        valeur finale.
\end{itemize}

% ======================================================================
\section{Écran d'accueil : avant de commencer}

\subsection{Saisie de l'identité}

Deux champs sont à remplir :
\begin{itemize}
  \item \textbf{NOM} : en majuscules (l'application le convertit
        automatiquement en majuscules à la frappe) ;
  \item \textbf{Prénom} : la première lettre de chaque mot est mise en
        majuscule automatiquement (par exemple \texttt{jean-pierre} devient
        \texttt{Jean-Pierre}, \texttt{marie-claire} devient \texttt{Marie-Claire}).
\end{itemize}

\begin{attention}
Les deux champs sont \textbf{obligatoires} (au moins deux caractères chacun).
Un message d'alerte apparaît si l'un des deux est vide ou trop court.
\end{attention}

\subsection{Option PAP (Projet d'Accueil Personnalisé)}

La case \emph{« Je dispose d'un PAP (tiers-temps accordé : 133 min 20 s) »} doit
être cochée \textbf{uniquement} si vous bénéficiez effectivement d'un PAP.

\begin{center}
\begin{tabular}{ll}
\toprule
\textbf{Situation} & \textbf{Durée de l'épreuve} \\
\midrule
Sans PAP  & \(100\) minutes exactement \\
Avec PAP  & \(100 \times \tfrac{4}{3} = 133\) min \(20\) s (soit \(2\,\)h\(13\)min\(20\)s) \\
\bottomrule
\end{tabular}
\end{center}

Le chronomètre en haut de l'écran tient compte automatiquement de ce choix.

\subsection{Contrainte des 100 minutes entre deux épreuves}

Pour éviter qu'un élève enchaîne plusieurs devoirs à la suite, l'application
mémorise dans le navigateur l'heure de départ du dernier devoir. Si moins de
\textbf{100 minutes} se sont écoulées, le bouton \emph{Commencer le devoir}
reste cliquable mais un message d'erreur affiche le temps restant à attendre
(au format \texttt{X min Y s}).

\begin{attention}
Cette contrainte est \textbf{active} dans la version courante. Elle protège
l'égalité de traitement entre les élèves et ne peut pas être contournée
depuis l'interface.
\end{attention}

\subsection{Lancement}

Cliquez sur \textbf{Commencer le devoir}. À partir de cet instant :
\begin{itemize}
  \item le chronomètre démarre (décompte affiché en haut de l'écran) ;
  \item la surveillance anti-fraude devient active (voir
        section~\ref{sec:surveillance}) ;
  \item la barre de progression en bas de page commence à se remplir.
\end{itemize}

% ======================================================================
\section{Écran d'examen : structure du sujet}

\subsection{Organisation}

Le sujet est présenté en \textbf{quatre parties}, chacune dans un cadre blanc
arrondi. Chaque partie contient :
\begin{itemize}
  \item un titre (\emph{Exercice 1 — \dots}, \emph{GROS PROBLÈME — \dots}) ;
  \item un court texte d'introduction avec les données (expression de \(u_n\),
        relation de récurrence, inégalités d'encadrement, \dots) ;
  \item une ou plusieurs questions numérotées \(1, 2, 3, \dots\)
\end{itemize}

\subsection{Les quatre parties et leurs exigences}

\begin{center}
\renewcommand{\arraystretch}{1.3}
\begin{tabular}{@{}p{3.1cm}p{7.8cm}p{3.5cm}@{}}
\toprule
\textbf{Partie} & \textbf{Contenu mathématique} & \textbf{Type de réponse} \\
\midrule
Exercice 1 &
Limites de polynômes (factorisation par le terme dominant) ; quotients de
polynômes (degrés égaux, degrés différents). &
Nombre, fraction, \(+\infty\) ou \(-\infty\). \\
\addlinespace
Exercice 2 &
Encadrement par \(-1 \leqslant \sin n \leqslant 1\), \(-1 \leqslant \cos n \leqslant 1\),
quotient de type \(\dfrac{a\,n + b\sin n}{c\,n + d\cos n}\). &
Fraction \(a/c\) ou \(+\infty/-\infty\). \\
\addlinespace
Exercice 3 &
Racines carrées : \(\sqrt{n^2 + a\,n + b} - n\), \(\sqrt{k^2 n^2 + a\,n + b} - k\,n\)
(expression conjuguée). &
Fraction \(a/2\) ou \(a/(2k)\). \\
\addlinespace
Gros problème &
Suite rationnelle \(u_n = \dfrac{a\,n+b}{c\,n+d}\) : forme indéterminée, limite,
écart \(v_n = u_n - \ell\), perturbation \(w_n = u_n + \tfrac{\sin n}{n}\),
produit \(z_n = n\,v_n\). &
Limite \(\ell\), \(0\), ou fraction. \\
\bottomrule
\end{tabular}
\end{center}

\subsection{Les lignes de raisonnement}

Sous chaque question se trouve une \textbf{zone de rédaction} composée de lignes.
La règle fondamentale est :

\begin{retenir}[Une idée = une ligne]
Chaque ligne doit contenir \textbf{une seule idée}, écrite en français, avec les
calculs nécessaires. On n'écrit pas un long paragraphe : on découpe le
raisonnement en étapes courtes, chacune sur sa propre ligne.
\end{retenir}

Le bouton \textbf{« + Ajouter une ligne de raisonnement »} ajoute autant de
lignes que nécessaire, à la fin de la liste.

\subsection{Les conjonctions de coordination}

\begin{center}
\fcolorbox{bleu}{blue!5}{%
\begin{minipage}{0.93\linewidth}
\textbf{Règle des conjonctions}
\begin{itemize}
  \item La \textbf{première ligne} d'une question ne commence par \textbf{aucun
        mot particulier} : la liste des conjonctions est désactivée
        (mention \emph{« début du raisonnement — pas de conjonction »}).
  \item Chaque ligne \textbf{suivante} doit commencer par \textbf{une} des sept
        conjonctions de coordination :
        \texttt{mais}, \texttt{ou}, \texttt{et}, \texttt{donc}, \texttt{or},
        \texttt{ni}, \texttt{car} — ou par le mot \texttt{Finalement}.
  \item On choisit la conjonction dans la \textbf{liste déroulante} située à
        gauche de la ligne, puis on écrit la suite du raisonnement dans la zone
        de texte.
\end{itemize}
\end{minipage}}
\end{center}

\subsection{La ligne « Finalement » : la réponse}

\begin{retenir}[La réponse est la ligne « Finalement »]
La \textbf{réponse} à une question est la ligne qui commence par le mot
\textbf{\texttt{Finalement}}. Dès que vous sélectionnez \texttt{Finalement} dans
la liste déroulante :
\begin{itemize}
  \item la ligne est mise en valeur (bordure orange épaisse, fond crème) ;
  \item le rendu du texte passe en gras dans l'aperçu ;
  \item c'est \emph{cette} ligne qui est analysée automatiquement pour la
        détection de la valeur finale (section~\ref{sec:detection}).
\end{itemize}
\end{retenir}

\begin{attention}[Conséquence à la remise de la copie]
Si une question ne comporte \textbf{aucune} ligne \texttt{Finalement} remplie,
le fichier \texttt{.json} indique \texttt{correcte: false} pour cette question :
le professeur verra que vous n'avez pas fourni de réponse formelle. Il est donc
essentiel de toujours terminer une question par une ligne \texttt{Finalement}.
\end{attention}

% ======================================================================
\section{Validation, aperçu et édition des lignes}
\label{sec:edition}

\subsection{La barre d'outils de chaque ligne}

Chaque ligne de raisonnement possède, à droite de son en-tête, une petite
\textbf{barre d'outils} discrète. Les boutons disponibles sont :

\begin{center}
\renewcommand{\arraystretch}{1.35}
\begin{tabular}{@{}c p{3.0cm} p{7.6cm}@{}}
\toprule
\textbf{Bouton} & \textbf{Nom} & \textbf{Action} \\
\midrule
\textbf{✓} & Valider & Analyse la ligne, insère les \texttt{\$} autour des
                     fragments mathématiques et affiche l'aperçu \LaTeX{} sous
                     la zone de saisie. \\
\addlinespace
\textbf{✎} & Modifier & Réapparaît uniquement sur les lignes validées :
                     déverrouille la ligne, efface l'aperçu, replace le curseur
                     dans la zone de saisie avec le texte sélectionné. \\
\addlinespace
\textbf{▲} & Monter & Échange la ligne avec celle du dessus. \\
\addlinespace
\textbf{▼} & Descendre & Échange la ligne avec celle du dessous. \\
\addlinespace
\textbf{✕} & Supprimer & Retire la ligne et son contenu (une confirmation est
                        demandée si du texte est présent). \\
\bottomrule
\end{tabular}
\end{center}

\begin{attention}
Les boutons \textbf{✎}, \textbf{▲}, \textbf{▼}, \textbf{✕} n'apparaissent que
lorsqu'ils ont un sens : par exemple, \textbf{▲} est masqué sur la première
ligne, \textbf{▼} sur la dernière, et \textbf{✕} lorsqu'il ne reste qu'une
seule ligne dans la question. La première ligne d'une question ne peut
jamais être supprimée.
\end{attention}

\subsection{Valider une ligne}

Il y a \textbf{trois} façons de valider une ligne :
\begin{enumerate}
  \item Appuyer sur \textbf{\texttt{Ctrl+Entrée}} (\texttt{⌘+Entrée} sur Mac)
        dans la zone de saisie ;
  \item Cliquer sur le bouton \textbf{✓} dans la barre d'outils de la ligne ;
  \item \textbf{Quitter le champ} : cliquer ailleurs sur la page, ou appuyer
        sur \texttt{Tab}. La validation se déclenche automatiquement lorsque le
        focus quitte la ligne.
\end{enumerate}

Après validation :
\begin{itemize}
  \item la bordure de la zone de saisie passe au \textbf{vert} ;
  \item un symbole \textbf{✓} apparaît après le numéro de la ligne ;
  \item un \textbf{aperçu \LaTeX{}} s'affiche sous la zone, avec la mention
        \emph{« Aperçu LaTeX : »} suivie de la phrase mise en forme.
\end{itemize}

\begin{retenir}[L'aperçu n'apparaît qu'après validation]
Contrairement à la frappe au kilomètre, l'aperçu ne se met \emph{pas} à jour à
chaque caractère. Il apparaît uniquement après la validation. Si vous modifiez
le texte ensuite, l'aperçu est effacé et le marquage \textbf{✓} disparaît :
il faut revalider.
\end{retenir}

\subsection{Modifier une ligne validée}

Pour reprendre une ligne déjà validée :
\begin{enumerate}
  \item Cliquez sur le bouton \textbf{✎} de cette ligne.
  \item La ligne redevient éditable : l'aperçu disparaît, la bordure repasse au
        gris, le curseur se place dans la zone de saisie et \emph{sélectionne}
        son contenu (vous pouvez donc réécrire directement par-dessus).
  \item Après votre correction, validez à nouveau (\texttt{Ctrl+Entrée} ou
        bouton \textbf{✓}).
\end{enumerate}

\subsection{Déplacer une ligne}

Les boutons \textbf{▲} et \textbf{▼} échangent la ligne avec sa voisine. Utile
pour réordonner un raisonnement sans le réécrire. Les aperçus déjà validés
sont \textbf{conservés} lors du déplacement.

\begin{attention}
Si vous déplacez une ligne vers la \textbf{première position}, sa conjonction
est automatiquement supprimée (la première ligne ne porte pas de conjonction).
Inversement, la ligne qui devient la deuxième reçoit la liste déroulante des
conjonctions, et vous devez en choisir une.
\end{attention}

\subsection{Supprimer une ligne}

Le bouton \textbf{✕} retire la ligne de la question. Une boîte de dialogue de
confirmation apparaît \emph{uniquement} si la ligne contient du texte. S'il ne
reste qu'une seule ligne, le bouton \textbf{✕} n'est pas affiché.

% ======================================================================
\section{Notations \LaTeX{} : \emph{sans} taper les \texttt{\$}}
\label{sec:notations}

\begin{retenir}[Nouveauté importante]
Vous n'avez \textbf{plus besoin} de taper les \texttt{\$} autour de vos
formules. Tapez le code \LaTeX{} directement :
l'application \emph{détecte} automatiquement les fragments mathématiques et
ajoute les \texttt{\$} à votre place.

Par exemple, si vous tapez :

\begin{center}
\verb|donc u_n = \dfrac{1}{n}|
\end{center}

l'application produit en interne :

\begin{center}
\verb|donc $u_n = \dfrac{1}{n}$|
\end{center}

puis affiche l'aperçu : \quad donc $u_n = \dfrac{1}{n}$.
\end{retenir}

\begin{attention}
Si vous tapez malgré tout des \texttt{\$} (ancien réflexe), l'application les
respecte et \emph{ne rajoute rien}. Les deux écritures sont donc compatibles :
avec ou sans \texttt{\$}, le rendu final est identique.
\end{attention}

\subsection{Comment fonctionne la détection automatique}

L'application découpe votre phrase en \emph{tokens} (mots et espaces), puis
classe chaque token :

\begin{itemize}
  \item \textbf{token mathématique fort} : il contient au moins un caractère
        typique de \LaTeX{} ou des mathématiques : \verb|\|, \verb|{|, \verb|}|,
        \verb|_|, \verb|^|, \verb|=|, \verb|*|, \verb|/|, \verb|<|, \verb|>|,
        ou bien les symboles $\infty$, $\leq$, $\geq$, $\neq$ ; ou encore c'est
        un nombre pur, ou une combinaison chiffre/opérateur ;
  \item \textbf{token mathématique doux} : une lettre isolée (par exemple
        \texttt{n}), un nom de fonction (\texttt{sin}, \texttt{cos},
        \texttt{lim}, \texttt{sqrt}, \dots), un nombre pur, une parenthèse
        seule, ou un identifiant suivi de parenthèses (\texttt{u(n)},
        \texttt{cos(n)}, \dots) ;
  \item \textbf{token texte} : tout le reste (mots français, articles,
        ponctuation).
\end{itemize}

Chaque token fort amorce un \emph{run} (une plage contiguë de maths), qui
s'étend ensuite à gauche et à droite à travers les tokens doux voisins. Les
runs ainsi formés sont encadrés par \texttt{\$...\$}.

\subsection{Écriture des indices et des exposants}

\begin{center}
\renewcommand{\arraystretch}{1.35}
\footnotesize
\begin{tabular}{@{}p{4.2cm}p{5.6cm}p{4.2cm}@{}}
\toprule
\textbf{Notation mathématique} & \textbf{À taper dans la zone (sans \texttt{\$})} & \textbf{Aperçu après validation} \\
\midrule
\(u_0\)      & \verb|u_0|               & \(u_0\) \\
\(u_1\)      & \verb|u_1|               & \(u_1\) \\
\(u_n\)      & \verb|u_n|               & \(u_n\) \\
\(u_{n+1}\)  & \verb|u_{n+1}|           & \(u_{n+1}\) \\
\(u_{n+2}\)  & \verb|u_{n+2}|           & \(u_{n+2}\) \\
\(n^2\)      & \verb|n^2|               & \(n^2\) \\
\(n^3\)      & \verb|n^3|               & \(n^3\) \\
\(2^n\)      & \verb|2^n|               & \(2^n\) \\
\(3^{n+1}\)  & \verb|3^{n+1}|           & \(3^{n+1}\) \\
\bottomrule
\end{tabular}
\end{center}

\begin{attention}[Accolades obligatoires]
Dès que l'indice ou l'exposant comporte plusieurs caractères, il faut
l'entourer d'\textbf{accolades} \verb|{...}| :
\begin{itemize}
  \item \verb|u_{n+1}| et non \verb|u_n+1| (qui donne \(u_n + 1\), différent) ;
  \item \verb|3^{n+1}| et non \verb|3^n+1| (qui donne \(3^n + 1\), différent) ;
  \item \verb|u_{n-1}| pour \(u_{n-1}\).
\end{itemize}
\end{attention}

\subsection{Fractions : \texttt{\textbackslash dfrac\{...\}\{...\}}}

\begin{center}
\renewcommand{\arraystretch}{1.35}
\footnotesize
\begin{tabular}{@{}p{4.2cm}p{6.0cm}p{4.0cm}@{}}
\toprule
\textbf{Notation mathématique} & \textbf{À taper (sans \texttt{\$})} & \textbf{Aperçu} \\
\midrule
\(\dfrac{3}{4}\)               & \verb|\dfrac{3}{4}|            & \(\dfrac{3}{4}\) \\
\(\dfrac{a}{c}\)              & \verb|\dfrac{a}{c}|            & \(\dfrac{a}{c}\) \\
\(\dfrac{a\,n+b}{c\,n+d}\)    & \verb|\dfrac{an+b}{cn+d}|      & \(\dfrac{an+b}{cn+d}\) \\
\(\dfrac{\sin n}{n}\)         & \verb|\dfrac{\sin n}{n}|       & \(\dfrac{\sin n}{n}\) \\
\(\dfrac{5}{n+1}\)            & \verb|\dfrac{5}{n+1}|          & \(\dfrac{5}{n+1}\) \\
\bottomrule
\end{tabular}
\end{center}

\begin{astuce}[Trois commandes équivalentes]
\verb|\dfrac|, \verb|\frac| et \verb|\tfrac| sont acceptées. Utilisez
\verb|\dfrac| (fraction « display ») pour qu'elle soit bien lisible dans
l'aperçu.
\end{astuce}

\subsection{Racines carrées : \texttt{\textbackslash sqrt\{...\}}}

\begin{center}
\renewcommand{\arraystretch}{1.35}
\footnotesize
\begin{tabular}{@{}p{4.2cm}p{6.0cm}p{4.0cm}@{}}
\toprule
\textbf{Notation mathématique} & \textbf{À taper (sans \texttt{\$})} & \textbf{Aperçu} \\
\midrule
\(\sqrt{n+1}\)          & \verb|\sqrt{n+1}|            & \(\sqrt{n+1}\) \\
\(\sqrt{n^2 + 3n + 5}\) & \verb|\sqrt{n^2+3n+5}|       & \(\sqrt{n^2+3n+5}\) \\
\(\sqrt{4n^2 - 5}\)     & \verb|\sqrt{4n^2-5}|         & \(\sqrt{4n^2-5}\) \\
\bottomrule
\end{tabular}
\end{center}

\subsection{Fonctions trigonométriques}

\begin{center}
\renewcommand{\arraystretch}{1.35}
\footnotesize
\begin{tabular}{@{}p{4.2cm}p{6.2cm}p{3.8cm}@{}}
\toprule
\textbf{Notation mathématique} & \textbf{À taper (sans \texttt{\$})} & \textbf{Aperçu} \\
\midrule
\(\sin n\)              & \verb|\sin n|                  & \(\sin n\) \\
\(\cos n\)              & \verb|\cos n|                  & \(\cos n\) \\
\(-1 \leqslant \sin n \leqslant 1\) &
\verb|-1 \leq \sin n \leq 1| & \(-1 \leq \sin n \leq 1\) \\
\bottomrule
\end{tabular}
\end{center}

\begin{astuce}[Commandes d'inégalité]
On peut écrire \verb|\leq| ou \verb|\le| pour \(\leq\), et \verb|\geq| ou
\verb|\ge| pour \(\geq\). Les versions \verb|\leqslant| et \verb|\geqslant|
sont aussi reconnues et donnent \(\leqslant\), \(\geqslant\).
\end{astuce}

\subsection{Limites : \texttt{\textbackslash lim}, \texttt{\textbackslash to}, \texttt{\textbackslash infty}}

\begin{center}
\renewcommand{\arraystretch}{1.35}
\footnotesize
\begin{tabular}{@{}p{4.8cm}p{6.8cm}p{3.0cm}@{}}
\toprule
\textbf{Notation mathématique} & \textbf{À taper (sans \texttt{\$})} & \textbf{Aperçu} \\
\midrule
\(\lim\limits_{n\to+\infty} u_n = +\infty\) &
\verb|\lim_{n \to +\infty} u_n = +\infty| & \(\lim\limits_{n\to+\infty} u_n = +\infty\) \\
\(\lim\limits_{n\to+\infty} \dfrac{1}{n} = 0\) &
\verb|\lim_{n \to +\infty} \dfrac{1}{n} = 0| & \(\lim\limits_{n\to+\infty} \dfrac{1}{n} = 0\) \\
\(+\infty\) & \verb|+\infty| & \(+\infty\) \\
\(-\infty\) & \verb|-\infty| & \(-\infty\) \\
\(n \to +\infty\) & \verb|n \to +\infty| & \(n \to +\infty\) \\
\bottomrule
\end{tabular}
\end{center}

\begin{attention}[Ne pas confondre \(\to\) et \(\mapsto\)]
Le signe « tend vers » s'écrit \verb|\to| (\(\to\)). Évitez \verb|\mapsto|
(\(\mapsto\)) qui signifie autre chose.
\end{attention}

\subsection{Parenthèses et multiplication}

\begin{center}
\renewcommand{\arraystretch}{1.35}
\footnotesize
\begin{tabular}{@{}p{4.6cm}p{6.6cm}p{3.8cm}@{}}
\toprule
\textbf{Notation mathématique} & \textbf{À taper (sans \texttt{\$})} & \textbf{Aperçu} \\
\midrule
\(n(n+1)\)            & \verb|n(n+1)|                & \(n(n+1)\) \\
\(2 \times u_{n+1}\)  & \verb|2 \times u_{n+1}|      & \(2 \times u_{n+1}\) \\
\((a\,n+b)(c\,n+d)\)  & \verb|(an+b)(cn+d)|          & \((an+b)(cn+d)\) \\
\(\left(\dfrac{1}{2}\right)^n\) &
\verb|\left(\dfrac{1}{2}\right)^n| & \(\left(\dfrac{1}{2}\right)^n\) \\
\bottomrule
\end{tabular}
\end{center}

\begin{astuce}[Le symbole de multiplication]
En \LaTeX, on utilise \verb|\times| pour \(\times\). Le point médian
\verb|\cdot| donne \(\cdot\). Les accolades \verb|\left( ... \right)| servent
à faire des parenthèses qui grandissent avec leur contenu.
\end{astuce}

\subsection{Exemples rédigés de bout en bout}

\subsubsection*{Exemple 1 — Polynôme de degré 3}

\emph{Énoncé type} : \og Pour tout \(n \in \mathbb{N}\), \(u_n = -2n^3 + 5n^2 - 7n + 3\).
Déterminer \(\lim\limits_{n \to +\infty} u_n\). \fg

\begin{center}
\footnotesize
\renewcommand{\arraystretch}{1.3}
\begin{tabular}{@{}p{2.4cm}p{11.8cm}@{}}
\toprule
\textbf{Conjonction} & \textbf{Contenu de la ligne (sans \texttt{\$})} \\
\midrule
\emph{(aucune)} & \verb|Je factorise par le terme dominant n^3| \\
\texttt{donc}   & \verb|u_n = n^3 (-2 + 5/n - 7/n^2 + 3/n^3)| \\
\texttt{or}     & \verb|lim_{n \to +\infty} 5/n = lim_{n \to +\infty} 7/n^2 = 0| \\
\texttt{donc}   & \verb|la parenthèse tend vers -2| \\
\texttt{or}     & \verb|lim_{n \to +\infty} n^3 = +\infty| \\
\texttt{donc}   & \verb|le produit -2 \times (+\infty) tend vers -\infty| \\
\texttt{Finalement} & \verb|lim_{n \to +\infty} u_n = -\infty| \\
\bottomrule
\end{tabular}
\end{center}

\subsubsection*{Exemple 2 — Quotient de mêmes degrés}

\emph{Énoncé type} : \og Pour tout \(n \geqslant 1\),
\(v_n = \dfrac{-3n^2 + 4n - 1}{2n^2 - n + 7}\). Déterminer
\(\lim\limits_{n \to +\infty} v_n\). \fg

\begin{center}
\footnotesize
\renewcommand{\arraystretch}{1.3}
\begin{tabular}{@{}p{2.4cm}p{11.8cm}@{}}
\toprule
\textbf{Conjonction} & \textbf{Contenu de la ligne (sans \texttt{\$})} \\
\midrule
\emph{(aucune)} & \verb|Je factorise numérateur et dénominateur par n^2| \\
\texttt{donc}   & \verb|v_n = \dfrac{n^2(-3 + 4/n - 1/n^2)}{n^2(2 - 1/n + 7/n^2)}| \\
\texttt{et}     & \verb|les n^2 se simplifient| \\
\texttt{donc}   & \verb|v_n = \dfrac{-3 + 4/n - 1/n^2}{2 - 1/n + 7/n^2}| \\
\texttt{or}     & \verb|les termes en 1/n et 1/n^2 tendent vers 0| \\
\texttt{donc}   & \verb|le quotient tend vers \dfrac{-3}{2}| \\
\texttt{Finalement} & \verb|lim_{n \to +\infty} v_n = \dfrac{-3}{2}| \\
\bottomrule
\end{tabular}
\end{center}

\subsubsection*{Exemple 3 — Racine carrée}

\emph{Énoncé type} : \og Pour tout \(n \in \mathbb{N}\),
\(w_n = \sqrt{n^2 + 6n + 5} - n\). Déterminer
\(\lim\limits_{n \to +\infty} w_n\). \fg

\begin{center}
\footnotesize
\renewcommand{\arraystretch}{1.3}
\begin{tabular}{@{}p{2.4cm}p{11.8cm}@{}}
\toprule
\textbf{Conjonction} & \textbf{Contenu de la ligne (sans \texttt{\$})} \\
\midrule
\emph{(aucune)} & \verb|Je multiplie par l'expression conjuguée| \\
\texttt{donc}   & \verb|w_n = \dfrac{(\sqrt{n^2+6n+5} - n)(\sqrt{n^2+6n+5} + n)}{\sqrt{n^2+6n+5} + n}| \\
\texttt{et}     & \verb|le numérateur se simplifie en n^2+6n+5 - n^2 = 6n+5| \\
\texttt{donc}   & \verb|w_n = \dfrac{6n+5}{\sqrt{n^2+6n+5} + n}| \\
\texttt{or}     & \verb|je factorise par n au numérateur et au dénominateur| \\
\texttt{donc}   & \verb|w_n = \dfrac{n(6 + 5/n)}{n(\sqrt{1 + 6/n + 5/n^2} + 1)}| \\
\texttt{Finalement} & \verb|lim_{n \to +\infty} w_n = 3| \\
\bottomrule
\end{tabular}
\end{center}

\subsubsection*{Exemple 4 — Encadrement par \(\sin\) et \(\cos\)}

\emph{Énoncé type} : \og Pour tout \(n \geqslant 1\),
\(t_n = \dfrac{3n + 2\sin n}{5n + 4\cos n}\). Déterminer
\(\lim\limits_{n \to +\infty} t_n\). \fg

\begin{center}
\footnotesize
\renewcommand{\arraystretch}{1.3}
\begin{tabular}{@{}p{2.4cm}p{11.8cm}@{}}
\toprule
\textbf{Conjonction} & \textbf{Contenu de la ligne (sans \texttt{\$})} \\
\midrule
\emph{(aucune)} & \verb|Pour tout n, -1 \leq \sin n \leq 1 et -1 \leq \cos n \leq 1| \\
\texttt{donc}   & \verb|-2 \leq 2\sin n \leq 2 et -4 \leq 4\cos n \leq 4| \\
\texttt{donc}   & \verb|3n - 2 \leq 3n + 2\sin n \leq 3n + 2| \\
\texttt{et}     & \verb|5n - 4 \leq 5n + 4\cos n \leq 5n + 4| \\
\texttt{or}     & \verb|je divise numérateur et dénominateur par n| \\
\texttt{donc}   & \verb|le quotient est encadré par des expressions tendant vers 3/5| \\
\texttt{Finalement} & \verb|lim_{n \to +\infty} t_n = \dfrac{3}{5}| \\
\bottomrule
\end{tabular}
\end{center}

% ======================================================================
\section{Détection automatique de la réponse}
\label{sec:detection}

\subsection{Comment la réponse est extraite}

L'application analyse votre ligne \texttt{Finalement} et extrait, dans l'ordre
d'apparition, \textbf{tous} les objets suivants :
\begin{enumerate}
  \item les fractions \verb|\dfrac{}{}|, \verb|\frac{}{}|, \verb|\tfrac{}{}|
        (les nombres sont convertis et réduits) ;
  \item les fractions numériques écrites avec une barre oblique
        (par exemple \verb|3/4|) ;
  \item les infinis : \verb|\infty|, \verb|\infinity|, \verb|infini|,
        \verb|infinie|, ou le symbole \(\infty\) lui-même, précédés
        éventuellement d'un signe ;
  \item les nombres (entiers ou décimaux, séparateur \texttt{.} ou \texttt{,}).
\end{enumerate}
La \textbf{dernière occurrence} trouvée dans la ligne est retenue comme
\textbf{réponse normalisée}, puis comparée à la réponse attendue.

\begin{retenir}[Règle d'or pour la ligne « Finalement »]
Terminez toujours votre ligne \texttt{Finalement} par la \textbf{valeur finale
de la limite}, clairement isolée. Par exemple :
\begin{itemize}
  \item \verb|Finalement lim_{n \to +\infty} u_n = +\infty| ;
  \item \verb|Finalement la limite est \dfrac{3}{2}| ;
  \item \verb|Finalement lim_{n \to +\infty} u_n = 0|.
\end{itemize}
Évitez les valeurs parasites après la conclusion (par exemple « \dots = 0,
car \(n\) est grand » : le mot \emph{grand} n'est pas un nombre, mais
d'autres nombres pourraient être captés à tort).
\end{retenir}

\subsection{Formes acceptées pour une même limite}

L'analyseur reconnaît indifféremment (que vous tapiez les \texttt{\$} ou non) :

\begin{center}
\footnotesize
\renewcommand{\arraystretch}{1.3}
\begin{tabular}{@{}p{3.4cm}p{8.6cm}p{2.8cm}@{}}
\toprule
\textbf{Valeur de la limite} & \textbf{Écritures acceptées} & \textbf{Résultat} \\
\midrule
\(\dfrac{3}{2}\)     & \verb|3/2|, \verb|\dfrac{3}{2}|, \verb|\frac{3}{2}|, \verb|1,5| & \texttt{3/2} \\
\(-\dfrac{5}{4}\)   & \verb|-5/4|, \verb|\dfrac{-5}{4}|, \verb|-1,25| & \texttt{-5/4} \\
\(0\)               & \verb|0| & \texttt{0} \\
\(+\infty\)         & \verb|+\infty|, \verb|+infini| & \texttt{+inf} \\
\(-\infty\)         & \verb|-\infty|, \verb|-infini| & \texttt{-inf} \\
\bottomrule
\end{tabular}
\end{center}

% ======================================================================
\section{Chronomètre, progression et surveillance}
\label{sec:surveillance}

\subsection{Chronomètre}

En haut de l'écran figure un chronomètre au format \texttt{MM:SS}. Il affiche le
temps \emph{restant}. La durée totale étant de \(100\) minutes, l'affichage
initial est \textbf{\texttt{100:00}} : il n'y a pas de bascule à \texttt{00:00}
au bout d'une heure, mais une descente continue jusqu'à zéro.

Lorsqu'il ne reste plus que \textbf{5 minutes} (300 secondes), le chronomètre
devient rouge et clignote pour vous alerter.

\subsection{Barre de progression}

La barre horizontale en bas de la page indique, en couleur, la chronologie de
la session :
\begin{itemize}
  \item \textbf{bleu} : temps de travail effectif sur la page ;
  \item \textbf{rouge} : périodes où la page est masquée ou la fenêtre
        minimisée (temps « hors page ») ;
  \item \textbf{orange} : périodes d'inactivité prolongée ;
  \item \textbf{violet} : instants où un \texttt{Ctrl+V} a été effectué dans
        une zone de réponse.
\end{itemize}
Les graduations sont espacées de \textbf{10 minutes} (adaptées à une épreuve
de 100 minutes).

\subsection{Surveillance}

L'application enregistre discrètement plusieurs indicateurs qui accompagneront
votre copie :
\begin{itemize}
  \item \textbf{Sorties de la page} : chaque changement d'onglet, minimisation
        ou perte de focus est compté, ainsi que la durée totale passée hors
        de la page.
  \item \textbf{Périodes d'inactivité} : toute plage sans clic, frappe ou
        défilement supérieure à \textbf{45 secondes} est enregistrée.
  \item \textbf{Collages} : chaque \texttt{Ctrl+V} dans une zone de réponse
        est compté, ainsi que le nombre de caractères collés.
  \item \textbf{Copies} : chaque \texttt{Ctrl+C} depuis une zone de saisie
        est compté.
  \item \textbf{Outils de développement} : la détection approximative d'une
        fenêtre d'outils ouverte est signalée.
\end{itemize}

Ces indicateurs alimentent un \textbf{taux estimé de fraude} (de \(0\) à
\(100~\%\)) calculé automatiquement et affiché au moment de la remise de la
copie, avec le détail des facteurs. Cet indice n'est \emph{pas} une preuve
de fraude : c'est une simple information destinée au professeur.

\subsection{Sauvegarde automatique}

Toutes les \textbf{15 secondes}, ainsi qu'à la fermeture de l'onglet, vos
réponses et le journal des événements sont sauvegardés dans la mémoire du
navigateur (\texttt{localStorage}). En cas de problème, une partie du travail
peut être récupérée.

% ======================================================================
\section{Remise de la copie}

\subsection{Le bouton « Terminer et enregistrer »}

Dans la barre supérieure se trouve un bouton vert \textbf{Terminer et
enregistrer}.

\begin{enumerate}
  \item L'application vérifie que chaque question comporte au moins une ligne
        \texttt{Finalement} remplie.
  \item Un message de confirmation vous est demandé (« Terminer le devoir
        maintenant ? »).
  \item Après confirmation, la copie est \textbf{définitivement figée} :
        le fichier \texttt{.json} est généré et téléchargé automatiquement,
        et la page tente de se fermer.
\end{enumerate}

\subsection{Fin du temps imparti}

Si le chronomètre arrive à zéro, l'application rend automatiquement la copie :
le fichier est téléchargé, une fenêtre récapitulative s'affiche, et la page
tente de se fermer.

\begin{attention}
Ne fermez pas la fenêtre vous-même avant la fin : la remise doit se faire par
le bouton \textbf{Terminer et enregistrer} (ou par la fin automatique du
temps). Une fermeture manuelle en cours d'épreuve déclenche un avertissement
du navigateur.
\end{attention}

% ======================================================================
\section{Le fichier \texttt{.json} produit}

\subsection{Nom du fichier}

Le fichier porte un nom de la forme :

\begin{center}
\verb|TG-MATH-COMPLEMENTAIRES_<NOM>-<Prenom>_<jjmmaaaa>.json|
\end{center}

Par exemple : \verb|TG-MATH-COMPLEMENTAIRES_DUPONT-Jean-Pierre_15102026.json|.

\subsection{Chiffrement}

Le contenu du fichier est \textbf{chiffré} :
\begin{itemize}
  \item \textbf{AES-256-GCM} avec dérivation de clé \textbf{PBKDF2-SHA256}
        (150\,000 itérations), lorsque le navigateur expose l'API
        \texttt{crypto.subtle} (cas normal sur un navigateur récent en
        HTTPS ou en \texttt{file://} sécurisé) ;
  \item un chiffrement de repli (XOR + Base64) est utilisé si cette API
        n'est pas disponible.
\end{itemize}
Le fichier reste lisible par le professeur qui dispose de la phrase de
passe ou de l'outil de déchiffrement fourni avec l'application.

\subsection{Contenu déchiffré}

Une fois déchiffré, le fichier contient l'intégralité de la copie :
\begin{itemize}
  \item l'identité de l'élève (nom, prénom, PAP) ;
  \item les horaires de début et de fin, la durée accordée et la durée
        réelle ;
  \item l'\textbf{état du délai entre essais} (\texttt{delaiEntreEssaisActif},
        \texttt{delaiEntreEssaisMs}, \texttt{delaiEntreEssaisMin}) ;
  \item la \textbf{durée de travail de base}
        (\texttt{dureeDevoirBaseMs}, \texttt{dureeDevoirBaseMin}) ;
  \item les \textbf{drapeaux de fonctionnalité} \texttt{autoWrapLatexActif}
        et \texttt{editionLignesActif} ;
  \item le \textbf{journal des événements} (sorties de page, inactivités,
        collages, ajouts, suppressions, déplacements, déverrouillages,
        validations de lignes, conjonctions choisies, ouverture de la notice
        de droits d'auteur, \dots) ;
  \item le \textbf{temps de travail par question} (en millisecondes) ;
  \item l'\textbf{énoncé généré} (avec les valeurs aléatoires réellement
        tirées) et la réponse attendue pour chaque question ;
  \item les \textbf{réponses} de l'élève, ligne par ligne, avec :
        \begin{itemize}
          \item la conjonction utilisée,
          \item le texte brut tapé (\texttt{texte}),
          \item la version avec \texttt{\$} ajoutés (\texttt{texteLatex}),
          \item l'état de validation (\texttt{validee}),
          \item le HTML de l'aperçu MathJax (\texttt{renduHtml}) ;
        \end{itemize}
        et pour la ligne \texttt{Finalement}, la \textbf{valeur normalisée
        détectée} et l'indication \texttt{correcte: true/false} ;
  \item l'\textbf{indice de fraude} détaillé (taux, niveau, facteurs).
\end{itemize}

% ======================================================================
\section{Dépannage / questions fréquentes}

\begin{description}[leftmargin=1.2em, style=nextline]
  \item[Le bouton « Commencer le devoir » affiche un message d'attente.]
        Un devoir a été commencé il y a moins de 100 minutes dans ce
        navigateur. Attendez le temps indiqué par le message.
  \item[Les formules de l'énoncé ne s'affichent pas.] Vérifiez votre
        connexion Internet : MathJax est chargé depuis un CDN. Sans
        connexion, les formules apparaissent en code \LaTeX{} brut.
  \item[Ma formule apparaît en texte brut dans l'aperçu.] Vous avez
        probablement oublié une partie du code \LaTeX{}. Vérifiez les
        accolades \verb|{...}|, les indices avec \verb|_| et les
        commandes avec \verb|\|.
  \item[Je ne vois pas d'aperçu sous ma zone de saisie.] L'aperçu
        n'apparaît qu'\textbf{après validation}. Appuyez sur
        \texttt{Ctrl+Entrée}, cliquez sur le bouton \textbf{✓}, ou cliquez
        ailleurs sur la page.
  \item[Je ne vois pas de réponse encadrée en orange.] Vous n'avez pas
        choisi \texttt{Finalement} dans la liste déroulante. Sélectionnez-le
        sur la ligne qui contient votre conclusion.
  \item[Ma valeur finale n'a pas été reconnue comme correcte alors que je
        pense avoir juste.] Vérifiez que la ligne \texttt{Finalement} se
        termine bien par la valeur numérique (nombre, fraction ou
        \(+\infty/-\infty\)). Relisez la section~\ref{sec:detection}.
  \item[Le bouton \textbf{✎} n'apparaît pas.] Ce bouton n'apparaît que sur
        les lignes déjà validées. Validez d'abord la ligne (\texttt{Ctrl+Entrée}
        ou \textbf{✓}), puis le bouton \textbf{✎} apparaîtra à côté.
  \item[Je ne peux pas supprimer une ligne.] Soit c'est la première ligne
        d'une question (qui doit toujours exister), soit c'est la seule
        ligne restante. Ajoutez d'abord une ligne avec le bouton
        \textbf{+ Ajouter une ligne de raisonnement}.
  \item[J'ai fermé la page par erreur.] Le travail est sauvegardé dans le
        navigateur toutes les 15 secondes, mais la copie n'est \emph{pas}
        transmise tant que le bouton \textbf{Terminer et enregistrer} n'a
        pas été utilisé.
  \item[Le fichier ne se télécharge pas.] Vérifiez que votre navigateur
        autorise les téléchargements depuis cette page, puis recommencez.
  \item[Que signifie la modale « © » ?] C'est la notice de droits
        d'auteur et de licence d'utilisation de l'application. Elle
        s'ouvre en cliquant sur le bouton rond \textbf{©} en bas à droite
        de l'écran, et se ferme par la touche \texttt{Échap}.
\end{description}

% ======================================================================
\section{Annexe : mémo à garder sous les yeux}

\begin{retenir}[Mémo \LaTeX{} pour les suites et limites]
\begin{center}
\footnotesize
\renewcommand{\arraystretch}{1.3}
\begin{tabular}{ll}
\toprule
\textbf{En mathématiques} & \textbf{À taper dans la réponse (sans \texttt{\$})} \\
\midrule
\(u_0\)                  & \verb|u_0| \\
\(u_1\)                  & \verb|u_1| \\
\(u_n\)                  & \verb|u_n| \\
\(u_{n+1}\)              & \verb|u_{n+1}| \\
\(u_{n+2}\)              & \verb|u_{n+2}| \\
\(n^2\)                  & \verb|n^2| \\
\(n^3\)                  & \verb|n^3| \\
\(2^n\)                  & \verb|2^n| \\
\(3^{n+1}\)              & \verb|3^{n+1}| \\
\(\dfrac{a}{b}\)         & \verb|\dfrac{a}{b}| \\
\(\dfrac{\sin n}{n}\)    & \verb|\dfrac{\sin n}{n}| \\
\(\sqrt{n^2+3n+5}\)      & \verb|\sqrt{n^2+3n+5}| \\
\(\sin n\), \(\cos n\)   & \verb|\sin n|, \verb|\cos n| \\
\(-1 \leq \sin n \leq 1\) & \verb|-1 \leq \sin n \leq 1| \\
\(+\infty\), \(-\infty\) & \verb|+\infty|, \verb|-\infty| \\
\(n \to +\infty\)        & \verb|n \to +\infty| \\
\(\lim\limits_{n\to+\infty} u_n\) & \verb|\lim_{n \to +\infty} u_n| \\
\(u_n = 0\)              & \verb|u_n = 0| \\
\(\dfrac{3}{2}\)         & \verb|\dfrac{3}{2}| \\
\bottomrule
\end{tabular}
\end{center}
\end{retenir}

\begin{retenir}[Mémo rédaction]
\begin{enumerate}
  \item \textbf{Une idée = une ligne.}
  \item \textbf{Première ligne} : aucune conjonction.
  \item \textbf{Lignes suivantes} : \texttt{mais, ou, et, donc, or, ni, car}
        ou \texttt{Finalement}.
  \item \textbf{Réponse} : la ligne qui commence par \texttt{Finalement}
        (encadrée en orange), terminée par la \textbf{valeur finale} de la
        limite.
  \item \textbf{Maths} : tapez le code \LaTeX{} directement
        (\verb|\dfrac|, \verb|\sqrt|, \verb|\infty|, \verb|\lim|,
        \verb|\to|), \textbf{sans \texttt{\$}}.
\end{enumerate}
\end{retenir}

\begin{retenir}[Mémo validation et édition]
\begin{itemize}
  \item \textbf{Valider} une ligne : \texttt{Ctrl+Entrée}, bouton
        \textbf{✓}, ou quitter le champ. L'aperçu apparaît alors.
  \item \textbf{Modifier} une ligne validée : bouton \textbf{✎}.
  \item \textbf{Déplacer} : boutons \textbf{▲} et \textbf{▼}.
  \item \textbf{Supprimer} : bouton \textbf{✕} (jamais sur la première ligne,
        ni sur la dernière restante).
\end{itemize}
\end{retenir}

\begin{retenir}[Mémo « ce qui est automatiquement détecté »]
L'analyseur extrait, dans l'ordre, la \textbf{dernière} :
\begin{itemize}
  \item fraction \verb|\dfrac{}{}|, \verb|\frac{}{}| ou \verb|\tfrac{}{}| ;
  \item fraction numérique \verb|a/b| ;
  \item mention d'infini : \verb|\infty|, \verb|\infinity|, \verb|infini|,
        ou le symbole \(\infty\), avec signe éventuel ;
  \item nombre entier ou décimal (\texttt{.} ou \texttt{,}).
\end{itemize}
D'où la règle : \emph{terminer la ligne \texttt{Finalement} par la valeur
finale}.
\end{retenir}

\vfill
\begin{center}
\small\itshape
Document rédigé par Yann Merdy — Licence d'utilisation non commerciale
(voir la modale « Droits d'auteur » dans l'application).
\end{center}

\end{document}